\subsection{More on Ideals}

\begin{prop}
Let \( I  \) and \( J  \) be ideals of a ring \( R  \). 
Define \( I + J  \) as \( \{ a + b : a \in I \ \text{and} \ b \in J  \}  \). Then \( I + J  \) is an ideal of \( R  \) and it is the smallest ideal that contains both \( I \) and \( J  \).
\end{prop}
\begin{proof}
Note that \( I  \) and \( J  \) are subrings of \( R  \) and so they are both normal additive subgroups of \( R  \). Hence, \( I + J  \) is an additive subgroup of \( R  \).

Let \( r (x + y) \in r (I + J ) \) where \( x \in I  \) and \( y \in J  \). Then 
\[  r(x + y) = rx + ry.  \]
But note that \( rx \in I  \) and \( ry \in J  \). Thus, \( r(x+y) \in I + J  \). Similarly, if \( (x+y)r \in (I + J)r  \), then 
\[  (x+y)r = xr + yr \]
where \( xr \in I  \) and \( yr \in J  \) since \( I  \) and \( J  \) are ideals of \( R  \). Thus, \( (x+y) r \in I + J  \). This tells us that \( I + J  \) is an ideal of \( R  \).

Notice that for containment, it follows that since \( I + J  \) is closed under addition, we have for every \( x \in I  \), \( x = x + {0}_{R} \) implies that \( x \in I +  J \) and likewise, if \( x \in J  \), then \( x = {0}_{R} + x  \) implies that \( x \in I + J  \). Thus, both ideals are contained in \( I + J  \). To see that \( I + J  \) is the smallest ideal containing \( I  \) and \( J  \), take any arbitrary ideal \( K  \) in \( R  \) such that \( K  \) contains both \( I  \) and \( J  \). Since \( K  \) is closed under addition, it follows that \( I + J  \subseteq K   \) and so \( I + J  \)
\end{proof}

Let \( R  \) be a ring with identity \( 1  \).

\begin{definition}[Ideal Generated By A Subset Of A Ring]
    Let \( A  \) be a subset of \( R  \). Let \( (A) \) denote the smallest ideal of \( R  \) containing \( A  \). We can define \( (A) \) more concretely by
    \[  (A) = \bigcap_{ I \in \Lambda }^{  }  I  \]
    where \( \Lambda = \{  \text{\( I \) is an ideal of \( R  \)} : A \subseteq I  \}  \). An ideal generated by a finite set is called a finitely generated ideal.
\end{definition}

If \( A = \{ {a}_{1}, {a}_{2}, \dots, {a}_{n} \}  \), then \( (A) \) will be written \( ({a}_{1}, {a}_{2}, \dots, {a}_{n}) \). Also, note that \( (I,J) = (I \cup J) \) (if the ideal is generated by two sets \( I  \) and \( J  \)).

What is the simplest group generated by a set? In group theory, we call these sets cyclic groups. They generate groups with single elements. In the ring setting, we call these kinds of sets \textbf{principal ideals}. That is, \( a \neq 0  \) in \( R  \) is a \textbf{principal ideal} if \( (a)  \) is an ideal. Here, \( (a) \) must be closed with respect to addition. More precisely, 

\begin{definition}[Principal Ideals]
    Let \( R  \) be a ring and \( a \in R  \). If \( R  \) is a commutative ring with \( 1 \neq 0  \), then the set
    \[  (a) = \{  ar : r \in R  \}  \]
    is called the \textbf{principal ideal} of \( R  \) generated by \( a  \).
\end{definition}

\begin{eg}
   \begin{enumerate}
       \item[(1)] In any ring, \( 0 = (0)  \) is a principal ideal generating the trivial ring and \( (1) = R  \) is the principal ideal that generates \( R  \).
       \item[(2)] In \( \Z  \), all ideals are \textbf{principal}. Why? All additive subgroups of \( \Z  \) are of the form \( n \Z  \) (\( n \geq 0  \)) and \( n \Z  \) is an ideal for all \( n \geq 0  \). Further, \( n\Z = \{  n x : x \in \Z  \}  = (n) \). Because \( \Z  \) only has principal ideals, we see that 
           \[  (6,9) = (d) = (3) \]
           where \( d = \text{gcd}(6,9) = 3   \).
           One can prove in general that for any \( a,b \in \Z  \) non-zero, \( (a,b) = (d) \) where \( d = \text{gcd}(a,b)  \).
   \end{enumerate} 
\end{eg}

\begin{prop}
   Let \( I  \) be an ideal of \( R  \) and let \( R  \) be a ring. 

   \begin{enumerate}
       \item[(1)] \( I  =R  \) if and only if \( I  \) contains a unit.
        \item[(2)] Assume \( R  \) is commutative. \( R  \) is a field if and only if the only ideals are \(  0  \) and \( R  \).
   \end{enumerate}
\end{prop}
\begin{proof}
\begin{enumerate}
    \item[(1)] If \( I = R  \), then \( 1 \in I  \). Conversely, if \( u  \) is a unit in \( I  \) with inverse \( v \in R  \), then for any \( r \in R  \), we have 
        \[  r = r \cdot 1 = r (vu) = (rv) u \in I.  \]
        Hence, \( R = I  \).
    \item[(2)] The ring \( R  \) is a field if and only if every non-zero element is a unit. Hence, any non-zero ideal contains a unit. Then the only nonzero ideal of \( R  \) is \( R  \) itself.

        Conversely, if \( 0  \) and \( R  \) are the only ideals of \( R  \), we let \( u \neq 0  \) in \( R  \) be arbitrary. By hypothesis, \( (u) = R  \) and so \( 1 \in (u) \). That is, 
        \begin{center}
            \( 1 = uv  \) for some \( v \in R  \) (also, \( vu = 1  \) since \( R  \) is commutative)
        \end{center}
        That is, \( u  \) is a unit. Thus, every nonzero element of \( R  \) is a unit so \( R  \) is a field.

\end{enumerate}
\end{proof}

\begin{corollary}
    If \( R  \) is a field, then any nonzero ring homomorphism from \( R  \) into another ring \( S  \) is an injection.
\end{corollary}
\begin{proof}
Let \( R  \) be a field. Let \( \varphi: R \to S  \) be a non-zero ring homomorphism. Since \( \text{ker}\varphi   \) is an ideal of the field \( R  \), \( \text{ker} \ \varphi = 0    \) or \( \text{ker} \varphi = R   \). If \( \text{ker} \ \varphi = R   \), then for all \( r \in R  \) \( \varphi(r) = 0  \) contrary to our assumption. Hence, \( \text{ker} \ \varphi = 0    \) and so \( \varphi  \) must be an injection.
\end{proof}

\begin{definition}[Maximal Ideal]
    An ideal \( M  \) in a ring \( S  \) is called a \textbf{maximal ideal} if it is proper (that is, \( M \neq S  \)) and the only ideals which contain \( M  \) are \( S  \) and itself; that is, if for all ideals \( K \leq R  \) with \( M \leq K \leq R  \), then \( K = M  \) or \( K = R  \).
\end{definition}

\begin{prop}
   In a ring with identity, every proper ideal is contained in a maximal ideal. (Be able to prove this for finite rings in the exam). 
\end{prop}

\begin{prop}[Proposition 12 from Dummit and Foote]
    Assume \( R  \) is commutative. The ideal \( M  \) is a maximal if and only if \( R / M  \) is a field.
\end{prop}
\begin{proof}
We see that \( M  \) is a maximal ideal if and only if there are no ideals, \( I  \), with \( M < I < R \). Using the fourth isomorphism theorem, the ideals of \( R  \) containing \( M  \) are in \( 1-1 \) correspondence with the ideals of \( R /M  \). This true if and only if the only ideals of \( R /M  \) are \( R /M  \) itself and \( \overline{0}  \). This true if and only if \( R / M  \) is a field.  
\end{proof}
\begin{remark}[Things to Verify In The Above Proof]
\begin{itemize}
    \item Show that if \( R   \) is commutative, then \( R / M  \) is commutative.
    \item \( R /M  \) has unity, and \( R /M  \) is not the zero ring. (This is true because \( M \neq R  \) which implies that \( R /M \neq \overline{0} \)).
\end{itemize}
\end{remark}

\begin{eg}
   What are the maximal ideals in \( \Z  \)?

   We know all subrings of \( \Z  \) are of the form in \( n\Z  \) for \( n \geq 0  \). It is straightforward to show that \( n \Z  \) is an ideal. 

   The maximal ideals are those ideals such that \( \Z / n \Z  \) is a field. We know that \( \Z / n \Z \) is a field if and only if \( n  \) is prime. That is , \( n \Z  \) is maximal if and only if \( p  \) is prime.
\end{eg}

\begin{definition}[Prime Ideals]
    Assume \( R  \) is a commutative ring. An ideal \( P  \) is a prime ideal if \( P \neq R  \) (this means that \( P  \) is not the zero ring) and if \( ab \in P  \), then either \( a \in P  \) or \( b\in P  \).
\end{definition}

\begin{prop}[Proposition 13 from Dummit and Foote]
    Assume \( R  \) is commutative. The ideal \( P  \) is a prime ideal in \( R  \) if and only if \( R /  P \) is an integral domain.
\end{prop}
\begin{proof}
The ideal \( P  \) is prime if and only if \( P \neq R  \) and if \( ab \in P  \), then either \( a \in P  \) or \( b \in P  \). That is, either \( a + P = {0}_{R} + P  \) or \( b + P  = {0}_{R} + P   \). Since \( P \neq R  \), we have \( R /P  \neq \overline{0} \) and if \( (a + P) (b+P) = 0 + P  \), then either \( a + P  = {0}_{R} + P  \) or \( b  + P  = {0}_{R} + P  \), implying that \( R /P  \) is an integral domain. 
\end{proof}


\begin{corollary}
    Every maximal ideal is a prime ideal when our ring commutative.
\end{corollary}

Suppose that \( M  \) is an ideal such that \( R /M  \) is a field, then \( M  \) is maximal.

\( 7 x^{13} - 11 x^{9} + 5 x^{5} - 2 x^{3} + 3  \) modulo \( x^{4} - 16 \)
